\Needspace{18\baselineskip}
\subsection{A countable divisorial positivity criterion}
\label{subsec:countable-divisorial-positivity}

A prescribed divisorial part can determine the ordinary type of an
ideal envelope even when it has infinitely many components.
On the projective plane with a pluriclosed background, the
smooth residual determines the mass of every representative of that type.
This gives a positivity test directly from divisorial orders.

Let $X=\mathbb P^2$, and let $\kappa$ be the Fubini--Study form
with $\int_X\kappa^2=1$. Choose distinct hyperplanes $H_j$ and
coefficients $a_j>0$ with $A=\sum_{j\geq1}a_j<\infty$.
For unit defining sections $s_j$, put
\[
 \ell_j(Z)=\log\frac{|s_j(Z)|^2}{\|Z\|^2}\leq0,
 \qquad u_0=\sum_{j\geq1}a_j\ell_j,
 \qquad u=u_0-\sup_Xu_0.
\]
Let $\theta$ be a smooth real $(1,1)$-form satisfying
$\ddc\theta=0$ and $R=\theta-A\kappa>0$.

\Needspace{14\baselineskip}
\begin{lemma}\label{countable-divisorial:profile}\index{Divisorial current!summable hyperplane series}
The series $u_0$ defines a quasi-plurisubharmonic function, and
\[
 \ddc u=S-A\kappa,\qquad
 S=\sum_{j\geq1}a_j[H_j],\qquad
 \theta+\ddc u=R+S>0.
\]
For every smooth compact modification $\pi\colon Z\to X$ and
every prime divisor $E$ on $Z$, the series defining the order
profile converges, and
\[
 b(E):=\sum_{j\geq1}a_j\operatorname{ord}_E(\pi^*H_j)
       =\nu(u\circ\pi,E).
\]
\end{lemma}
\begin{proof}
The Poincar\'e--Lelong formula gives
$\ddc\ell_j=[H_j]-\kappa$. By unitary invariance, the integral
of $\ell_j$ against the probability measure $\kappa^2$ is
computed from the distribution of $x=|Z_0|^2/\|Z\|^2$, whose
density on $[0,1]$ is $2(1-x)$. Thus
\[
 \int_X\ell_j\kappa^2=\int_0^1 2(1-x)\log x\,\mathrm dx=-\frac32.
\]
Since $\ell_j\leq0$, the partial sums of $u_0$ decrease, and
they are $A\kappa$-plurisubharmonic. By monotone convergence,
$\int_Xu_0\kappa^2=-\frac32A$ is finite, so the limit $u_0$ is
not identically $-\infty$ and is $A\kappa$-plurisubharmonic.
The partial sums converge in $L^1$, so their complex Hessians
converge as currents; this gives the current identities, and
$\theta+\ddc u\geq R>0$.

Fix $\pi$ and $E$, and choose a Gauduchon form $\gamma$ on the
compact surface $Z$. Consider the tail
$r_N=\sum_{j>N}a_j\ell_j$ of total coefficient
$A_N=\sum_{j>N}a_j$. Its pullback $r_N\circ\pi$ is
$A_N\pi^*\kappa$-plurisubharmonic; it is not identically
$-\infty$ because $\pi$ is biholomorphic over a dense open
subset of $X$. Hence the closed positive current
\[
 T_N=A_N\pi^*\kappa+\ddc(r_N\circ\pi)
\]
has trace mass $\int_ZT_N\wedge\gamma=A_N\int_Z\pi^*\kappa\wedge\gamma$,
by distributional Stokes and $\ddc\gamma=0$.
By \cref{lem:nonclosed-divisorial-removal},
$T_N\geq\nu(r_N\circ\pi,E)[E]$, and therefore
\[
 0\leq\nu(r_N\circ\pi,E)
 \leq A_N\frac{\int_Z\pi^*\kappa\wedge\gamma}{\int_E\gamma}.
\]
The denominator is positive, so this bound tends to zero.
Lelong numbers are additive, and
$\nu(\ell_j\circ\pi,E)=\operatorname{ord}_E(\pi^*H_j)$.
Hence the $N$th partial sum of $b(E)$ differs from
$\nu(u\circ\pi,E)$ by $\nu(r_N\circ\pi,E)$.
This proves convergence and the asserted equality.
\end{proof}

\Needspace{16\baselineskip}
\begin{proposition}\label{countable-divisorial:positive-mass}\index{Multiplier ideal!positive singularity volume}\index{Projective space!countable divisorial ideal type}
Let $v\in\PSH(X,\theta)$. The conditions
\[
 \nu(v\circ\pi,E)=b(E)\text{ for every }\pi,E,
 \qquad v\sim_{\mathcal I}u
\]
are equivalent. Under either condition,
\[
 q=P_\theta[v]_{\mathcal I}=P_\theta[u]_{\mathcal I}\sim u,
\]
and every $w\in\PSH(X,\theta)$ with $w\sim q$ satisfies
\begin{equation}\label{eq:countable-divisorial-positive-mass}
 \int_X\mathrm{MA}_\theta(w)
 =\underline{\operatorname{vol}}_\theta(q)
 =\overline{\operatorname{vol}}_\theta(q)
 =\int_XR^2>0.
\end{equation}
Moreover, $\theta$ is big and $X$ has the bounded mass property.
\end{proposition}
\begin{proof}
\emph{Ideal profile and envelope.} Choose an integer $C$
with $C\kappa\geq\theta$. Then $C>A$, since $\theta>A\kappa$,
and $u$ and $v$ are plurisubharmonic for the K\"ahler form
$C\kappa$. In this common closed background, the stated
profile implies $v\sim_{\mathcal I}u$ by the projective valuative
characterization of ideal equivalence, whose source is given in
\cref{sec:history-volumes}. On the surface $X$ this also follows
from the first item of \cref{thm:local-i-inputs}, which reduces
ideal equivalence to equality of the Lelong numbers at all points
of all local modifications. Over a smaller neighborhood, a local
modification of a surface is a finite composition of point blowups,
and these extend to a smooth compact modification of $X$. One further
point blowup turns the Lelong number at a point into the generic
Lelong number along the exceptional curve. Thus the profile
determines all these Lelong numbers.
Conversely, if $v\sim_{\mathcal I}u$, then
$v\circ\pi\sim_{\mathcal I}u\circ\pi$ for every smooth compact
modification $\pi$ by \cref{lem:birational-ideal-order}. The first
item of \cref{thm:local-i-inputs}, applied with the identity of $Z$,
then gives equal Lelong numbers at all points of $Z$, and hence
equal generic Lelong numbers along every prime divisor $E$.

The potential $u$ satisfies $\ddc u=S-A\kappa$, and $\theta=R+A\kappa$
is a Hermitian form. Hence \cref{cor:divisorial-envelope-equality},
applied with the smooth closed form $A\kappa$ and the Hermitian form
$\theta$, shows that the ordinary and ideal envelopes of $u$ agree
and are ordinarily equivalent to $u$.
Since $v\sim_{\mathcal I}u$, the families defining
$P_\theta[v]_{\mathcal I}$ and $P_\theta[u]_{\mathcal I}$ coincide,
which proves the assertion about $q$. The inequality
$\theta+\ddc u\geq R>0$ proves
bigness. For any Hermitian form $h$ on $X$, choose $D$ with
$D\kappa\geq h$. Bounded $h$-plurisubharmonic potentials are
$D\kappa$-plurisubharmonic, so positivity of bounded mixed
products and Stokes' theorem bound their masses by $\int_X(D\kappa)^2=D^2$.
This proves the bounded mass property.

\emph{The bounded residual.} Fix $w\in\PSH(X,\theta)$ with
$w\sim u$, and put $T=C\kappa+\ddc w$. Ordinary equivalence
preserves generic Lelong numbers, so \cref{countable-divisorial:profile},
applied with $\pi$ the identity, gives $\nu(w,H_j)=\nu(u,H_j)=a_j$.
\Cref{lem:nonclosed-divisorial-removal}, applied to the finite partial
sums of $S$, and passage to the limit give
\[
 K=T-S=(C-A)\kappa+\ddc g\geq0,\qquad g=w-u.
\]
Since $w\sim u$, the function $g$ is bounded off the poles of $u$.
By the inequality just displayed, it agrees almost everywhere with a
quasi-plurisubharmonic function, which is bounded by
\cref{lem:qpsh-almost-everywhere-order}. Thus $K$
has bounded local potentials, and its trace charges no
pluripolar set. The current $S$ is carried by the pole
union $\bigcup_jH_j\subset\{w=-\infty\}$. Thus
$\langle T\rangle=K$. For every smooth real $(1,1)$-form $L$
with $\ddc L=0$, distributional Stokes gives
\[
 \int_X L\wedge\langle T\rangle
       =(C-A)\int_XL\wedge\kappa.
\]

\emph{Closed top mass and smooth correction.} Put
$T_u=C\kappa+\ddc u=S+(C-A)\kappa$.
For the finite partial sums $u_N=\sum_{j\leq N}a_j\ell_j$,
the non-pluripolar mass of $(C\kappa+\ddc u_N)^2$ is
$(C-\sum_{j\leq N}a_j)^2$: outside $\bigcup_{j\leq N}H_j$
the current is the smooth form $(C-\sum_{j\leq N}a_j)\kappa$, and the
non-pluripolar product gives that union zero mass. Since $u_0\leq u_N$,
closed-background mass monotonicity, the closed case of
\cref{thm:volume-monotonicity}, gives
\[
 \int_X\langle T_u^2\rangle
 \leq\left(C-\sum_{j\leq N}a_j\right)^2
\]
for every $N$, and hence $\int_X\langle T_u^2\rangle\leq(C-A)^2$.
For the reverse inequality, locally write $S=\ddc\psi$ with
$\psi$ plurisubharmonic. The currents
$(C-A)\kappa+\ddc\max(\psi,-k)$ have bounded local potentials
and are at least $(C-A)\kappa$.
On the plurifine open set $\{\psi>-k\}$, their squared products
coincide with $\langle T_u^2\rangle$ and therefore
dominate $(C-A)^2\kappa^2$. Letting $k$ increase, and noting
that $\kappa^2$ does not charge the pluripolar set $\{\psi=-\infty\}$, gives
$\langle T_u^2\rangle\geq(C-A)^2\kappa^2$.
The two bounds give $\int_X\langle T_u^2\rangle=(C-A)^2$.
Since $C\kappa$ is a K\"ahler form and $w-u$ is bounded,
invariance of closed masses under bounded changes then gives
\[
 \int_X\langle T^2\rangle=(C-A)^2.
\]

Set $L=C\kappa-\theta\geq0$, so $\ddc L=0$.
Expanding the bounded products of $\max(w,-k)$ in the two
backgrounds and restricting to $\{w>-k\}$ gives, in the limit,
\[
 \mathrm{MA}_\theta(w)
       =\langle T^2\rangle-2L\wedge\langle T\rangle+L^2.
\]
All three measures on the right are finite: the first
has the mass just computed, $\langle T\rangle=K$, and
the smooth measure $L^2$ gives the pole set zero mass.
The same truncations are used in both backgrounds, so this
identity is an equality of Borel measures.
Integrating, and using the Stokes identity for $\int_XL\wedge\langle T\rangle$, yields
\[
 \int_X\mathrm{MA}_\theta(w)
 =(C-A)^2-2(C-A)\int_XL\wedge\kappa+\int_XL^2
 =\int_X(\theta-A\kappa)^2.
\]
This holds throughout the ordinary type, proving both
singularity-volume identities and their strict positivity.
\end{proof}

\Needspace{16\baselineskip}
\begin{example}\label{countable-divisorial:nonclosed-example}\index{Pluriclosed metric!nonanalytic positive-volume type}
There is a nonclosed pluriclosed background on $\mathbb P^2$
to which \cref{countable-divisorial:positive-mass} applies with a
nonanalytic ideal type: no potential with analytic singularities is
ideal equivalent to $u$. The prescribed
hyperplanes are not contained in a single pencil, and both
singularity volumes of the ideal model $q$ are at least $1/16$.
\end{example}
\begin{proof}
In a coordinate ball choose a smooth real cutoff $\chi$ equal
to one near the origin, and extend
$\alpha=\chi|z_2|^2\mathrm dz_1$ by zero. Put
\[
 \tau=\partial\overline\alpha+\bar\partial\alpha.
\]
It is real and smooth, with $\ddc\tau=0$. Where $\chi=1$,
\[
 \mathrm d\tau=\mathrm dz_2\wedge\mathrm d\bar z_2
                  \wedge(\mathrm dz_1-\mathrm d\bar z_1)\ne0.
\]
Choose $\varepsilon>0$ small enough that
$\theta=\kappa+\varepsilon\tau\geq\kappa/2$; then $\ddc\theta=0$
and $\mathrm d\theta=\varepsilon\,\mathrm d\tau\neq0$.
Set $A=1/4$, $a_j=A2^{-j}$, $t_j=2^{-j}$ and
\[
 H_j=\{Z_0+t_jZ_1+t_j^2Z_2=0\}.
\]
The Vandermonde determinant shows that any three defining
linear forms are independent; hence the hyperplanes are not
in one pencil. Here $R=\theta-A\kappa\geq\kappa/4$, so
$\int_XR^2\geq1/16$. \Cref{countable-divisorial:positive-mass} gives the
mass throughout the ordinary type and both singularity volumes.
By its first condition with $\pi$ the identity, every potential
ideal equivalent to $u$ has generic Lelong number $a_j>0$ along
each of the infinitely many lines $H_j$. A potential with analytic
singularities on a compact manifold has positive generic Lelong
numbers only along the finitely many prime divisors in its pole set.
Hence no analytic potential has this ideal type.
\end{proof}
